\subsection{Verification and measurement}\label{sec:verification}

The source carries 19 test suites, run on every build: the configuration parser, the
ring buffers, the JPEG decoding and framing, resolution negotiation, the operators and
the chain, the visibility correction, the camera controls, the raw recorder and the
webcam output, child processes, the frame bus under concurrent processes (including a
writer overtaking a reader during a copy), the health records, the HTTP server and
client, \hellion{}'s pipelines, relay, status file and RTP round trip through FFmpeg, the
H.264 packetizer packet by packet, and the resource measurement described below. A
camera's own H.264 was tested from a file source through the packetizer, the topside's
relay and decoder, and the recorder: every decoded frame and every frame of the
recording was byte-identical to FFmpeg's decode of the file. It was also run through
V4L2 from a loopback device carrying H.264 at 1280$\times$720 and 30~frames per
second, but not yet on a camera that encodes. On every change the suites run with GCC
and Clang, with and without OpenMP, on x86-64, and natively on AArch64, which executes
the NEON paths, and the HTTP parsers are fuzzed with libFuzzer for a minute. The
suites also run, as built for release, on ARMv7 and AArch64 under QEMU user-mode
emulation, and the suites of the changed components were run under AddressSanitizer
and UndefinedBehaviorSanitizer.

Each change that should not alter the output was checked against the previous code by
hashing the output of both over a fixed set of inputs: three 1280$\times$720 camera
photographs and synthetic frames of edge-case sizes (odd chroma widths, frames one or
two pixels wide), every strength and threshold of the operator in steps, 1 to 7
threads, on x86-64 and on both ARM targets under emulation (Table~\ref{tab:equivalence}).
Processing costs were measured as the minimum over repeated runs on an AMD Ryzen~5
5600 (6 cores) with the threads of OpenMP set to sleep while idle, as \eeye{}
configures them; costs on ARM were estimated as executed instruction counts under QEMU
with an instruction-counting plugin, which reflects the work but not the memory
behaviour of a Raspberry Pi (Table~\ref{tab:costs}).

The cost of a running installation is measured from outside by a command of
\hellion{} that follows every \eeye{} and \hellion{} process on the machine and the
child processes they started. Once a second it samples, from \texttt{/proc} and
\texttt{/sys} only, the processor time of each process and of each named thread,
resident memory, context switches, the processor's temperature, clock and throttling
state, the frame rates on the frame bus, the frame rate, drops, bit rate and loss of
each pipeline, and the traffic of each network interface; it does not touch the
processes it measures. On a Raspberry Pi~5, which has no hardware H.264 encoder, a
1280$\times$720 stream in placement~B took 0.91 of a core in all, of which the FFmpeg
encoder took 0.51; the Pi was throttled by heat during that run, so both figures are
upper estimates. This is the cost that sending a camera's own H.264
(Section~\ref{sec:transport}) removes from the vehicle; the saving itself is still to
be measured.

Transport latency was measured from frame bus to frame bus through two \hellion{}
processes, one as the vehicle and one as the topside, connected over loopback on an
idle laptop (AMD Ryzen~5 5600H, FFmpeg~9.0.2). A test program published frames whose
pixels encode their frame number and read them from the receiving bus; the latency of a
frame is the difference between the two times, and an RTP tap between the encoder and
the decoder divided it into stages (Table~\ref{tab:latency}). The complete chain was
also run between two machines on a local network, the laptop's built-in camera (HP
Wide Vision HD, UVC) serving as the vehicle's camera, for placements A and B, for the
camera controls in placement~A, for the recording and the decoder across disconnections
of the network cable, and for automatic pairing. Neither program has yet been run on the
vehicle, with its camera or across its tether.

\begin{table*}[!t]
  \centering\small
  \caption{Transport latency, frame bus to frame bus through two \hellion{} processes
    over loopback; median (10th--90th percentile) of 300 frames.}\label{tab:latency}
  \begin{adjustbox}{max width=\linewidth}
  \begin{tabular}{@{}lccc@{}}
    \toprule
    & \makecell{640$\times$480\\30 frames/s} & \makecell{640$\times$480\\60 frames/s}
    & \makecell{1280$\times$720\\30 frames/s} \\
    \midrule
    Without the delimiter and slice threading & 70.4 (70.2--70.6)~ms & 37.0 (36.8--37.2)~ms & 74.8 (74.6--75.1)~ms \\
    With both & 3.8 (3.7--4.1)~ms & 3.7 (3.6--3.9)~ms & 8.3 (8.1--8.6)~ms \\
    \bottomrule
  \end{tabular}
  \end{adjustbox}
  \tablenote{Without the two measures the decoder held two frame periods: about one in
    FFmpeg's H.264 parser and one in its raw video output. The encoder, from the frame
    written to its last RTP packet, took 1.2--1.5~ms at 640$\times$480, and \hellion{}'s
    own relay, pipe and bus about 0.5~ms. Camera exposure, the tether and the vehicle's
    encoder add to these figures.}
\end{table*}

\begin{table*}[!t]
  \centering\small
  \caption{Processing costs per 1280$\times$720 frame. x86-64: AMD Ryzen~5 5600,
    minimum over repeated runs, 4 threads unless stated. ARM: executed instructions under
    emulation, one thread, 20 frames.}\label{tab:costs}
  \setlength{\tabcolsep}{4pt}
  \begin{tabular}{@{}L{6.2cm}L{3.2cm}L{4.2cm}@{}}
    \toprule
    Stage & x86-64 & ARM (instructions) \\
    \midrule
    Image statistics and visibility correction & 0.55~ms (2.09~ms, 1 thread) & -- \\
    Blur, strength 1 & 0.18~ms & AArch64 183\,M; ARMv7 263\,M \\
    Blur, strength 4 & 0.59--0.66~ms & AArch64 283\,M; ARMv7 391\,M \\
    Blur, strength 8 & 0.84~ms & AArch64 486\,M; ARMv7 676\,M \\
    Blur, strength 20 & 1.56~ms & -- \\
    Canny, strength 2, thresholds 40/90 & 3.41--3.82~ms & -- \\
    Frame bus publication & 34.9~$\mu$s; 17.2~$\mu$s with neutral chroma & -- \\
    \addlinespace
    \hellion{} preview, 1 and 4 viewers & 4\% and 4\% of a core & -- \\
    Encoder (FFmpeg, libx264), 30 frames/s & 0.21~s of CPU per second of video, laptop & -- \\
    \bottomrule
  \end{tabular}
  \tablenote{Before the measures described in the text, the visibility correction took
    1.49~ms, the blur 0.75, 1.38 and 2.27~ms at strengths 1, 4 and 8 (and 900\,M and
    1\,108\,M instructions on AArch64 at 4 and 8), Canny 3.80--4.20~ms, the preview
    13\% of a core with four viewers, and the encoder 0.30~s of CPU per second of
    video.}
\end{table*}

\begin{table*}[!t]
  \centering\small
  \caption{Equivalence checks: output hashed over all cases for the previous and the
    revised code; every comparison was identical.}\label{tab:equivalence}
  \setlength{\tabcolsep}{4pt}
  \begin{tabular}{@{}L{3.6cm}L{6.2cm}rL{3.2cm}@{}}
    \toprule
    Change & Inputs & Cases & Platforms (threads) \\
    \midrule
    Visibility correction & smooth, flat and noisy synthetic sequences of 60 frames, at
      sizes from 1920$\times$1080 to 1$\times$1, including padded strides & -- & x86-64, AArch64, ARMv7 \\
    Canny & 3 photographs $\times$ 4 strengths $\times$ 7 low $\times$ 6 high thresholds & 504 & x86-64 (1, 3, 4), AArch64, ARMv7 (4) \\
    Blur, half-resolution NEON & 4 sizes $\times$ strengths 0--20 $\times$ 3 or 1 planes & 168 & x86-64 (1, 4), AArch64, ARMv7 (4) \\
    Blur, scratch-plane passes & 7 sizes, including 2-pixel-wide frames $\times$ strengths 0--20 $\times$ 3 or 1 planes & 294 & x86-64 (1, 3, 4), AArch64, ARMv7 (4) \\
    \bottomrule
  \end{tabular}
\end{table*}
